Table~\ref{tab:main} gives test accuracy (mean $\pm$ std over five seeds) and Fig.~\ref{fig:bars} plots it; Table~\ref{tab:stats} gives paired differences of the SimCLR-style probe against each baseline.

\begin{table}[t]\centering\caption{Main results: test accuracy, mean $\pm$ std over 5 seeds (best bold, second underlined).}\label{tab:main}
\resizebox{0.9\columnwidth}{!}{\begin{tabular}{llcc}
\toprule
Method & Task & test\_acc $\uparrow$ & $n$ \\
\midrule
Pixels + LR & n50 & 0.831 $\pm$ 0.040 & 5 \\
Random CNN + probe & n50 & 0.856 $\pm$ 0.022 & 5 \\
Supervised scratch & n50 & \underline{0.918 $\pm$ 0.017} & 5 \\
Rotation (reimplemented) & n50 & 0.772 $\pm$ 0.037 & 5 \\
PCA + LR & n50 & 0.820 $\pm$ 0.031 & 5 \\
\textbf{SimCLR-style probe} & n50 & \textbf{0.955 $\pm$ 0.012} & 5 \\
\midrule
Pixels + LR & n10 & 0.564 $\pm$ 0.061 & 5 \\
Random CNN + probe & n10 & 0.546 $\pm$ 0.118 & 5 \\
Supervised scratch & n10 & \underline{0.662 $\pm$ 0.043} & 5 \\
Rotation (reimplemented) & n10 & 0.441 $\pm$ 0.033 & 5 \\
PCA + LR & n10 & 0.505 $\pm$ 0.063 & 5 \\
\textbf{SimCLR-style probe} & n10 & \textbf{0.826 $\pm$ 0.037} & 5 \\
\midrule
Pixels + LR & n200 & 0.932 $\pm$ 0.013 & 5 \\
Random CNN + probe & n200 & 0.944 $\pm$ 0.015 & 5 \\
Supervised scratch & n200 & \underline{0.976 $\pm$ 0.010} & 5 \\
Rotation (reimplemented) & n200 & 0.908 $\pm$ 0.010 & 5 \\
PCA + LR & n200 & 0.914 $\pm$ 0.010 & 5 \\
\textbf{SimCLR-style probe} & n200 & \textbf{0.979 $\pm$ 0.010} & 5 \\
\bottomrule
\end{tabular}
}\end{table}

\begin{figure}[t]\centering\includegraphics[width=\columnwidth]{bars_main_test_acc.pdf}
\caption{Test accuracy by system and label budget (mean over 5 seeds).}\label{fig:bars}\end{figure}

\begin{table}[t]\centering\caption{SimCLR-style probe minus baseline: mean paired difference in accuracy, paired $t$-test $p$, Cohen's $d$ (5 seeds).}\label{tab:stats}
\resizebox{0.9\columnwidth}{!}{\begin{tabular}{llrrr}
\toprule
Task & Baseline & $\Delta$ acc & paired $p$ & $d$ \\
\midrule
n10 & PCA + LR & \rhval{cmp/main/pca-+-lr/n10/test_acc/delta:3} & \rhval{cmp/main/pca-+-lr/n10/test_acc/paired_p:4} & \rhval{cmp/main/pca-+-lr/n10/test_acc/cohen_d:1} \\
n10 & Pixels + LR & \rhval{cmp/main/pixels-+-lr/n10/test_acc/delta:3} & \rhval{cmp/main/pixels-+-lr/n10/test_acc/paired_p:4} & \rhval{cmp/main/pixels-+-lr/n10/test_acc/cohen_d:1} \\
n10 & Random CNN + probe & \rhval{cmp/main/random-cnn-+-probe/n10/test_acc/delta:3} & \rhval{cmp/main/random-cnn-+-probe/n10/test_acc/paired_p:4} & \rhval{cmp/main/random-cnn-+-probe/n10/test_acc/cohen_d:1} \\
n10 & Rotation (reimplemented) & \rhval{cmp/main/rotation-reimplemented/n10/test_acc/delta:3} & \rhval{cmp/main/rotation-reimplemented/n10/test_acc/paired_p:4} & \rhval{cmp/main/rotation-reimplemented/n10/test_acc/cohen_d:1} \\
n10 & Supervised scratch & \rhval{cmp/main/supervised-scratch/n10/test_acc/delta:3} & \rhval{cmp/main/supervised-scratch/n10/test_acc/paired_p:4} & \rhval{cmp/main/supervised-scratch/n10/test_acc/cohen_d:1} \\
\midrule
n50 & PCA + LR & \rhval{cmp/main/pca-+-lr/n50/test_acc/delta:3} & \rhval{cmp/main/pca-+-lr/n50/test_acc/paired_p:4} & \rhval{cmp/main/pca-+-lr/n50/test_acc/cohen_d:1} \\
n50 & Pixels + LR & \rhval{cmp/main/pixels-+-lr/n50/test_acc/delta:3} & \rhval{cmp/main/pixels-+-lr/n50/test_acc/paired_p:4} & \rhval{cmp/main/pixels-+-lr/n50/test_acc/cohen_d:1} \\
n50 & Random CNN + probe & \rhval{cmp/main/random-cnn-+-probe/n50/test_acc/delta:3} & \rhval{cmp/main/random-cnn-+-probe/n50/test_acc/paired_p:4} & \rhval{cmp/main/random-cnn-+-probe/n50/test_acc/cohen_d:1} \\
n50 & Rotation (reimplemented) & \rhval{cmp/main/rotation-reimplemented/n50/test_acc/delta:3} & \rhval{cmp/main/rotation-reimplemented/n50/test_acc/paired_p:4} & \rhval{cmp/main/rotation-reimplemented/n50/test_acc/cohen_d:1} \\
n50 & Supervised scratch & \rhval{cmp/main/supervised-scratch/n50/test_acc/delta:3} & \rhval{cmp/main/supervised-scratch/n50/test_acc/paired_p:4} & \rhval{cmp/main/supervised-scratch/n50/test_acc/cohen_d:1} \\
\midrule
n200 & PCA + LR & \rhval{cmp/main/pca-+-lr/n200/test_acc/delta:3} & \rhval{cmp/main/pca-+-lr/n200/test_acc/paired_p:4} & \rhval{cmp/main/pca-+-lr/n200/test_acc/cohen_d:1} \\
n200 & Pixels + LR & \rhval{cmp/main/pixels-+-lr/n200/test_acc/delta:3} & \rhval{cmp/main/pixels-+-lr/n200/test_acc/paired_p:4} & \rhval{cmp/main/pixels-+-lr/n200/test_acc/cohen_d:1} \\
n200 & Random CNN + probe & \rhval{cmp/main/random-cnn-+-probe/n200/test_acc/delta:3} & \rhval{cmp/main/random-cnn-+-probe/n200/test_acc/paired_p:4} & \rhval{cmp/main/random-cnn-+-probe/n200/test_acc/cohen_d:1} \\
n200 & Rotation (reimplemented) & \rhval{cmp/main/rotation-reimplemented/n200/test_acc/delta:3} & \rhval{cmp/main/rotation-reimplemented/n200/test_acc/paired_p:4} & \rhval{cmp/main/rotation-reimplemented/n200/test_acc/cohen_d:1} \\
n200 & Supervised scratch & \rhval{cmp/main/supervised-scratch/n200/test_acc/delta:3} & \rhval{cmp/main/supervised-scratch/n200/test_acc/paired_p:4} & \rhval{cmp/main/supervised-scratch/n200/test_acc/cohen_d:1} \\
\bottomrule
\end{tabular}
}\end{table}

\begin{figure}[t]\centering\includegraphics[width=0.8\columnwidth]{learning_curves.pdf}
\caption{Label-efficiency curves (mean $\pm$ std over 5 seeds). SimCLR is highest at every budget and the spread between systems narrows as labels grow.}\label{fig:curves}\end{figure}

\textbf{H1 (SimCLR $>$ scratch): supported at 10 and 50 labels, not at 200.} The paired differences are positive at every budget, but at 200 labels the gain over scratch is 0.002 with $p=0.18$ (Table~\ref{tab:stats}), which is within noise; with five seeds we cannot claim an advantage there. The advantage grows as labels shrink (0.164 at $n=10$, 0.037 at $n=50$).

\textbf{H2 (SimCLR $>$ rotation): supported} at all budgets, with large differences (Table~\ref{tab:stats}). The result is stronger than the hypothesis anticipated: rotation is below the untrained random encoder at all three budgets (clearly at 50 and 200 labels, $p=0.02$ and $0.004$; at 10 labels $p=0.09$) and not better than PCA + LR at any budget: below it at 10 and 50 labels (paired $p=0.05$ and $0.10$, not resolved with five seeds) and about equal at 200 (Tables~\ref{tab:main},~\ref{tab:extra}). The rotation pretraining loss in the registry is near zero, so the pretext task is solved easily (digits have an unambiguous upright orientation) and may yield features specialised to orientation (a hypothesis we did not test). We did not test remedies (harder pretext tasks, fewer epochs) and do not know whether a better-tuned rotation baseline would do better; rotation here is a failure of this configuration.

\textbf{H3 (SSL vs PCA): partly refuted.} SimCLR beats PCA + LR at 10 labels by 0.321, but the gap at 200 labels is still 0.065, far from the two-point band, so the ``gap shrinks to within 2 points'' part is refuted. The gap does shrink (0.321, 0.135, 0.065). Rotation pretraining does not beat PCA at 10 labels (0.441 vs 0.505), so the first part fails for rotation.

\begin{table}[t]\centering\caption{Per-seed test accuracy at 10 labels (seed = split, labeled subset and initialisation).}\label{tab:perseed}
\resizebox{0.9\columnwidth}{!}{\begin{tabular}{lrrrrr}
\toprule
System & s0 & s1 & s2 & s3 & s4 \\
\midrule
Supervised scratch & 0.608 & 0.624 & 0.696 & 0.700 & 0.682 \\
Random CNN + probe & 0.386 & 0.543 & 0.720 & 0.537 & 0.545 \\
SimCLR-style probe & 0.805 & 0.785 & 0.881 & 0.817 & 0.843 \\
\bottomrule
\end{tabular}
}\end{table}

\textbf{Where the method is unreliable.} At 10 labels all systems have large seed variance, notably the random encoder ($\pm0.118$). Table~\ref{tab:perseed} shows that SimCLR is above scratch and above the random encoder on each of the five seeds at $n=10$, but each seed has only one labeled example per class, so the labeled-subset draw dominates the noise. Part of the SimCLR advantage over pixels and PCA is attributable to the CNN architecture rather than pretraining: the random frozen encoder beats PCA at 50 and 200 labels and pixels at 200 labels (Table~\ref{tab:extra}), but not clearly pixels at 50 ($p=0.12$) or either at 10, where its variance is large.

\begin{table*}[t]\centering\footnotesize\caption{Additional paired-by-seed comparisons from the run registry (A minus B, 5 seeds; ``A wins'' counts seeds). $p$-values are uncorrected: 24 tests are shown here and 15 in Table~\ref{tab:stats}, and the three budgets share the same pretrained encoder and test split per seed, so they are not independent replications.}\label{tab:extra}
\begin{tabular}{llrr}
\toprule
Task & A $-$ B & $\Delta$ acc & paired $p$ \\
\midrule
n10 & SimCLR-style probe $-$ SimCLR geom-only & \rhval{cmp/cmp_aug/simclr-geom-only/n10/test_acc/delta:3} & \rhval{cmp/cmp_aug/simclr-geom-only/n10/test_acc/paired_p:3} \\
n50 & SimCLR-style probe $-$ SimCLR geom-only & \rhval{cmp/cmp_aug/simclr-geom-only/n50/test_acc/delta:3} & \rhval{cmp/cmp_aug/simclr-geom-only/n50/test_acc/paired_p:3} \\
n200 & SimCLR-style probe $-$ SimCLR geom-only & \rhval{cmp/cmp_aug/simclr-geom-only/n200/test_acc/delta:3} & \rhval{cmp/cmp_aug/simclr-geom-only/n200/test_acc/paired_p:3} \\
\midrule
n10 & SimCLR-style probe $-$ SimCLR photo-only & \rhval{cmp/cmp_aug/simclr-photo-only/n10/test_acc/delta:3} & \rhval{cmp/cmp_aug/simclr-photo-only/n10/test_acc/paired_p:3} \\
n50 & SimCLR-style probe $-$ SimCLR photo-only & \rhval{cmp/cmp_aug/simclr-photo-only/n50/test_acc/delta:3} & \rhval{cmp/cmp_aug/simclr-photo-only/n50/test_acc/paired_p:3} \\
n200 & SimCLR-style probe $-$ SimCLR photo-only & \rhval{cmp/cmp_aug/simclr-photo-only/n200/test_acc/delta:3} & \rhval{cmp/cmp_aug/simclr-photo-only/n200/test_acc/paired_p:3} \\
\midrule
n10 & Supervised + aug $-$ Supervised scratch & \rhval{cmp/cmp_supaug/supervised-scratch/n10/test_acc/delta:3} & \rhval{cmp/cmp_supaug/supervised-scratch/n10/test_acc/paired_p:3} \\
n50 & Supervised + aug $-$ Supervised scratch & \rhval{cmp/cmp_supaug/supervised-scratch/n50/test_acc/delta:3} & \rhval{cmp/cmp_supaug/supervised-scratch/n50/test_acc/paired_p:3} \\
n200 & Supervised + aug $-$ Supervised scratch & \rhval{cmp/cmp_supaug/supervised-scratch/n200/test_acc/delta:3} & \rhval{cmp/cmp_supaug/supervised-scratch/n200/test_acc/paired_p:3} \\
\midrule
n10 & Supervised + aug $-$ SimCLR-style probe & \rhval{cmp/cmp_supaug/simclr-style-probe/n10/test_acc/delta:3} & \rhval{cmp/cmp_supaug/simclr-style-probe/n10/test_acc/paired_p:3} \\
n50 & Supervised + aug $-$ SimCLR-style probe & \rhval{cmp/cmp_supaug/simclr-style-probe/n50/test_acc/delta:3} & \rhval{cmp/cmp_supaug/simclr-style-probe/n50/test_acc/paired_p:3} \\
n200 & Supervised + aug $-$ SimCLR-style probe & \rhval{cmp/cmp_supaug/simclr-style-probe/n200/test_acc/delta:3} & \rhval{cmp/cmp_supaug/simclr-style-probe/n200/test_acc/paired_p:3} \\
\midrule
n10 & Rotation (reimplemented) $-$ PCA + LR & \rhval{cmp/cmp_rotation/pca-+-lr/n10/test_acc/delta:3} & \rhval{cmp/cmp_rotation/pca-+-lr/n10/test_acc/paired_p:3} \\
n50 & Rotation (reimplemented) $-$ PCA + LR & \rhval{cmp/cmp_rotation/pca-+-lr/n50/test_acc/delta:3} & \rhval{cmp/cmp_rotation/pca-+-lr/n50/test_acc/paired_p:3} \\
n200 & Rotation (reimplemented) $-$ PCA + LR & \rhval{cmp/cmp_rotation/pca-+-lr/n200/test_acc/delta:3} & \rhval{cmp/cmp_rotation/pca-+-lr/n200/test_acc/paired_p:3} \\
\midrule
n10 & Rotation (reimplemented) $-$ Random CNN + probe & \rhval{cmp/cmp_rotation/random-cnn-+-probe/n10/test_acc/delta:3} & \rhval{cmp/cmp_rotation/random-cnn-+-probe/n10/test_acc/paired_p:3} \\
n50 & Rotation (reimplemented) $-$ Random CNN + probe & \rhval{cmp/cmp_rotation/random-cnn-+-probe/n50/test_acc/delta:3} & \rhval{cmp/cmp_rotation/random-cnn-+-probe/n50/test_acc/paired_p:3} \\
n200 & Rotation (reimplemented) $-$ Random CNN + probe & \rhval{cmp/cmp_rotation/random-cnn-+-probe/n200/test_acc/delta:3} & \rhval{cmp/cmp_rotation/random-cnn-+-probe/n200/test_acc/paired_p:3} \\
\midrule
n10 & Random CNN + probe $-$ PCA + LR & \rhval{cmp/cmp_random/pca-+-lr/n10/test_acc/delta:3} & \rhval{cmp/cmp_random/pca-+-lr/n10/test_acc/paired_p:3} \\
n50 & Random CNN + probe $-$ PCA + LR & \rhval{cmp/cmp_random/pca-+-lr/n50/test_acc/delta:3} & \rhval{cmp/cmp_random/pca-+-lr/n50/test_acc/paired_p:3} \\
n200 & Random CNN + probe $-$ PCA + LR & \rhval{cmp/cmp_random/pca-+-lr/n200/test_acc/delta:3} & \rhval{cmp/cmp_random/pca-+-lr/n200/test_acc/paired_p:3} \\
\midrule
n10 & Random CNN + probe $-$ Pixels + LR & \rhval{cmp/cmp_random/pixels-+-lr/n10/test_acc/delta:3} & \rhval{cmp/cmp_random/pixels-+-lr/n10/test_acc/paired_p:3} \\
n50 & Random CNN + probe $-$ Pixels + LR & \rhval{cmp/cmp_random/pixels-+-lr/n50/test_acc/delta:3} & \rhval{cmp/cmp_random/pixels-+-lr/n50/test_acc/paired_p:3} \\
n200 & Random CNN + probe $-$ Pixels + LR & \rhval{cmp/cmp_random/pixels-+-lr/n200/test_acc/delta:3} & \rhval{cmp/cmp_random/pixels-+-lr/n200/test_acc/paired_p:3} \\
\bottomrule
\end{tabular}
\end{table*}
